\subsection{Funtores de hipertor global de complejos de módulos
cuasicoherentes y sucesiones espectrales de Künneth: caso general}
\label{subsection:III.6.7}

\begin{env}[6.7.1]
\label{III.6.7.1}
Vamos a generalizar ahora las definiciones de (6.6.8) al caso en que $S$ es
un preesquema cualquiera, $X^{(i)}$ e $Y^{(i)}$ son $S$-preesquemas y
$f_i:X^{(i)}\to Y^{(i)}$ son morfismos \emph{separados y cuasicompactos}. Se
trata entonces, para todo par de complejos $\mathscr P_\bullet^{(i)}$ de
$\mathscr O_{X^{(i)}}$-módulos cuasicoherentes, acotados inferiormente
($i=1,2$), de definir para todo $n$ un $\mathscr O_Y$-módulo cuasicoherente
$\mathfrak{Tor}_n^S
(f_1,f_2;\mathscr P_\bullet^{(1)},\mathscr P_\bullet^{(2)})$, así como seis
funtores espectrales, que se reduzcan a las definiciones de (6.6.8) cuando
$S$, $Y^{(1)}$ e $Y^{(2)}$ son afines (se pone
$Y=Y^{(1)}\times_S Y^{(2)}$).

Supongamos primero que $S=\operatorname{Spec}(A)$ es \emph{afín}, pero que
$Y^{(1)}$ e $Y^{(2)}$ son cualesquiera; sea $W^{(i)}$ un abierto afín de
$Y^{(i)}$; entonces $f_i^{-1}(W^{(i)})$ es un $S$-esquema cuasicompacto y
$W=W^{(1)}\times_S W^{(2)}$ es un abierto afín de $Y$; sea
$f'_i:f_i^{-1}(W^{(i)})\to W^{(i)}$ la restricción de $f_i$, y sea
$\mathscr P'_\bullet{}^{(i)}$ la restricción
$\mathscr P_\bullet^{(i)}|f_i^{-1}(W^{(i)})$ ($i=1,2$). Se tienen entonces,
por (6.6.8), las sucesiones espectrales
${}^{(t)}\mathscr E
(f'_1,f'_2;\mathscr P'_\bullet{}^{(1)},\mathscr P'_\bullet{}^{(2)})$
de $(\mathscr O_Y|W)$-módulos cuasicoherentes, y hay que verificar que
satisfacen las condiciones de pegado ($0_{\mathrm I}$, 3.3.1).
Nos reducimos inmediatamente al caso en que
$Y^{(i)}=\operatorname{Spec}(B_i)$ es afín y $W^{(i)}=D(g_i)$, donde
$g_i\in B_i$, de modo que
$W=D(g_1\otimes g_2)$
\oldpage[III]{158}
en $Y=\operatorname{Spec}(B_1\otimes_A B_2)$ (II, 4.3.2.1); si
$X'^{(i)}=f_i^{-1}(W^{(i)})$, se trata de establecer un isomorfismo canónico
de funtores espectrales
\begin{env}[6.7.1.1]
\label{III.6.7.1.1}
\[
  {}^{(t)}\mathscr E
  (X'^{(1)},X'^{(2)};
   \mathscr P'_\bullet{}^{(1)},\mathscr P'_\bullet{}^{(2)})
  \xrightarrow{\sim}
  {}^{(t)}\mathscr E
  (X^{(1)},X^{(2)};
   \mathscr P_\bullet^{(1)},\mathscr P_\bullet^{(2)})
  \otimes_B B_g
  \tag{6.7.1.1}\label{III.6.7.1.1.eq}
\]
\end{env}

donde se ha puesto $B=B_1\otimes_A B_2$ y $g=g_1\otimes g_2$. Para ello,
partamos de cubrimientos abiertos afines finitos $\mathfrak U^{(i)}$ de
$X^{(i)}$ ($i=1,2$), y sea $\mathfrak U'^{(i)}$ la traza de
$\mathfrak U^{(i)}$ sobre $X'^{(i)}$, que sigue estando formada por abiertos
afines (I, 5.5.10); más precisamente, se tiene
\[
  C^\bullet(\mathfrak U'^{(i)},\mathscr P'_\bullet{}^{(i)})
  =C^\bullet(\mathfrak U^{(i)},\mathscr P_\bullet^{(i)})
   \otimes_{B_i}(B_i)_{g_i}.
\]
Si se pone
$L'_{\bullet\bullet}{}^{(i)}
=C^\bullet(\mathfrak U'^{(i)},\mathscr P'_\bullet{}^{(i)})$, se tiene, por
tanto,
\[
  L'_{\bullet\bullet}{}^{(1)}\otimes_A L'_{\bullet\bullet}{}^{(2)}
  =\bigl(L_{\bullet\bullet}^{(1)}\otimes_{B_1}(B_1)_{g_1}\bigr)
   \otimes_A
   \bigl(L_{\bullet\bullet}^{(2)}\otimes_{B_2}(B_2)_{g_2}\bigr);
\]
como $B_g=(B_1)_{g_1}\otimes_A(B_2)_{g_2}$ salvo isomorfismo canónico, se
tiene, también salvo isomorfismo canónico,
\[
  L'_{\bullet\bullet}{}^{(1)}\otimes_A L'_{\bullet\bullet}{}^{(2)}
  =(L_{\bullet\bullet}^{(1)}\otimes_A L_{\bullet\bullet}^{(2)})
    \otimes_B B_g.
\]
Si $M_{\bullet\bullet}^{(i)}$ es una resolución proyectiva de
Cartan--Eilenberg de $L_{\bullet\bullet}^{(i)}$, que puede suponerse formada
por $B_i$-módulos, del hecho de que $(B_i)_{g_i}$ es \emph{plano} sobre $B_i$
se deduce que
$M'_{\bullet\bullet}{}^{(i)}
=M_{\bullet\bullet}^{(i)}\otimes_{B_i}(B_i)_{g_i}$ es una resolución
proyectiva de Cartan--Eilenberg del bicomplejo
$L'_{\bullet\bullet}{}^{(i)}$; además, se tiene
\[
  M'_{\bullet\bullet}{}^{(1)}\otimes_A M'_{\bullet\bullet}{}^{(2)}
  =(M_{\bullet\bullet}^{(1)}\otimes_A M_{\bullet\bullet}^{(2)})
    \otimes_B B_g.
\]
El isomorfismo buscado (6.7.1.1) se deduce entonces inmediatamente de las
definiciones de la hiperhomología de un bicomplejo y de la exactitud del
funtor $G\otimes_B B_g$ respecto del $B$-módulo $G$.
\end{env}

\begin{env}[6.7.2]
\label{III.6.7.2}
Supongamos ahora que $S$ es cualquiera, y sean
$u_i:Y^{(i)}\to S$ los morfismos estructurales ($i=1,2$). Sea $(S_\alpha)$ un
cubrimiento abierto afín de $S$, pongamos
$Y_\alpha^{(i)}=u_i^{-1}(S_\alpha)$,
$X_\alpha^{(i)}=f_i^{-1}(Y_\alpha^{(i)})$, y sea
$f_{i\alpha}:X_\alpha^{(i)}\to Y_\alpha^{(i)}$ la restricción de $f_i$, que
es un morfismo \emph{separado y cuasicompacto}. Los
$Y_\alpha=Y_\alpha^{(1)}\times_{S_\alpha}Y_\alpha^{(2)}$ forman un
cubrimiento abierto de $Y$, y sobre cada $Y_\alpha$ quedan definidos, por
(6.7.1), los funtores espectrales
\[
  {}^{(t)}\mathscr E^{S_\alpha}
  \bigl(f_{1\alpha},f_{2\alpha};
    \mathscr P_\bullet^{(1)}|X_\alpha^{(1)},
    \mathscr P_\bullet^{(2)}|X_\alpha^{(2)}\bigr);
\]
hay que demostrar de nuevo que estos funtores verifican las condiciones de
pegado.

Nos reducimos inmediatamente a la situación siguiente:
$S=\operatorname{Spec}(A)$ es afín, $S'=D(h)$, donde $h\in A$, y
$u_i(Y^{(i)})\subset S'$; puede suponerse además que
$Y^{(i)}=\operatorname{Spec}(B_i)$ es afín; se trata de definir isomorfismos
canónicos
\begin{env}[6.7.2.1]
\label{III.6.7.2.1}
\[
  {}^{(t)}\mathscr E^S
  (X^{(1)},X^{(2)};
   \mathscr P_\bullet^{(1)},\mathscr P_\bullet^{(2)})
  \xrightarrow{\sim}
  {}^{(t)}\mathscr E^{S'}
  (X^{(1)},X^{(2)};
   \mathscr P_\bullet^{(1)},\mathscr P_\bullet^{(2)}).
  \tag{6.7.2.1}\label{III.6.7.2.1.eq}
\]
\end{env}
Ahora bien, con las notaciones de (6.6.2), los
$L_{\bullet\bullet}^{(i)}$ están formados por $A_h$-módulos, y se tiene
$L_{\bullet\bullet}^{(1)}\otimes_{A_h}L_{\bullet\bullet}^{(2)}
=L_{\bullet\bullet}^{(1)}\otimes_A L_{\bullet\bullet}^{(2)}$ salvo
isomorfismo canónico; como puede tomarse una resolución proyectiva de
Cartan--Eilenberg $M_{\bullet\bullet}^{(i)}$ de
$L_{\bullet\bullet}^{(i)}$ formada por $A_h$-módulos, esto proporciona
inmediatamente el isomorfismo canónico buscado.
\end{env}

Hemos demostrado, en resumen, el siguiente

\begin{theorem}[6.7.3]
\label{III.6.7.3}
Sean $S$ un preesquema, $f_i:X^{(i)}\to Y^{(i)}$ un $S$-morfismo separado y
cuasicompacto de $S$-preesquemas y $\mathscr P_\bullet^{(i)}$ un complejo
acotado inferiormente de $\mathscr O_{X^{(i)}}$-módulos cuasicoherentes
($i=1,2$); se pone $Y=Y^{(1)}\times_S Y^{(2)}$. Existe un
bi-$\partial$-funtor
  $\mathfrak{Tor}_\bullet^S
(f_1,f_2;\mathscr P_\bullet^{(1)},\mathscr P_\bullet^{(2)})$
con valores en la categoría de los $\mathscr O_Y$-módulos cuasicoherentes,
tal que, si $V^{(i)}$ es un abierto afín de $Y^{(i)}$ ($i=1,2$) y
$V=V^{(1)}\times_S V^{(2)}$, se tiene
\[
  \mathfrak{Tor}_\bullet^S
  (f_1,f_2;\mathscr P_\bullet^{(1)},\mathscr P_\bullet^{(2)})|V
  =\Bigl(\mathbf{Tor}_\bullet^S
    \bigl(f_1^{-1}(V^{(1)}),f_2^{-1}(V^{(2)});
      \mathscr P_\bullet^{(1)}|f_1^{-1}(V^{(1)}),
      \mathscr P_\bullet^{(2)}|f_2^{-1}(V^{(2)})\bigr)\Bigr)^{\sim}.
\]
\oldpage[III]{159}
Este bifuntor es el término límite de seis funtores espectrales
birregulares
\[
  {}^{(t)}\mathscr E
  (f_1,f_2;\mathscr P_\bullet^{(1)},\mathscr P_\bullet^{(2)})
  \qquad(t=a,b,a',b',c,d)
\]
cuyos términos $E_2$ vienen dados por
\[
\begin{aligned}
  {}^{(a)}\mathscr E_{pq}^2
  &=\bigoplus_{q_1+q_2=q}
    \mathscr{T}\!or_p^S
    \bigl(\mathscr H^{-q_1}(f_1,\mathscr P_\bullet^{(1)}),
          \mathscr H^{-q_2}(f_2,\mathscr P_\bullet^{(2)})\bigr),\\
  {}^{(b)}\mathscr E_{pq}^2
  &=\mathscr H^{-p}\bigl(f_1\times_S f_2,
    \mathfrak{Tor}_q^S
    (\mathscr P_\bullet^{(1)},\mathscr P_\bullet^{(2)})\bigr),\\
  {}^{(a')}\mathscr E_{pq}^2
  &=\bigoplus_{q_1+q_2=q}
    \mathfrak{Tor}_p^S
    \bigl(\mathscr H^{-q_1}(f_1,\mathscr P_\bullet^{(1)}),
          \mathscr H^{-q_2}(f_2,\mathscr P_\bullet^{(2)})\bigr),\\
  {}^{(b')}\mathscr E_{pq}^2
  &=\mathscr H^{-p}\bigl(f_1\times_S f_2,
    \mathscr{T}\!or_q^S
    (\mathscr P_\bullet^{(1)},\mathscr P_\bullet^{(2)})\bigr),\\
  {}^{(c)}\mathscr E_{pq}^2
  &=\bigoplus_{q_1+q_2=q}
    \mathfrak{Tor}_p^S
    \bigl(f_1,f_2;
      \mathscr H_{q_1}(\mathscr P_\bullet^{(1)}),
      \mathscr H_{q_2}(\mathscr P_\bullet^{(2)})\bigr),\\
  {}^{(d)}\mathscr E_{pq}^2
  &=\mathscr H_p\bigl(\mathfrak{Tor}_q^S
    (f_1,f_2;\mathscr P_\bullet^{(1)},\mathscr P_\bullet^{(2)})\bigr).
\end{aligned}
\]
Se dice que las sucesiones espectrales $(a)$ y $(b)$ son las
\emph{sucesiones espectrales de Künneth}.

Obsérvese que las sucesiones espectrales $(a)$ y $(a')$ (respectivamente,
$(b)$ y $(b')$) son idénticas cuando $\mathscr P_\bullet^{(1)}$ y
$\mathscr P_\bullet^{(2)}$ se reducen a sus términos de grado 0; en este caso,
las sucesiones $(c)$ y $(d)$ son degeneradas y, por tanto, no presentan
interés.
\end{theorem}

\begin{remark}[6.7.4]
\label{III.6.7.4}
Los hipertor globales que acabamos de definir comprenden como casos
particulares tanto los \emph{módulos de hipercohomología} definidos en
(6.2.1) como los \emph{hipertor locales} definidos en (6.5.3). Demostremos que
para todo morfismo $f:X\to Y$ \emph{cuasicompacto y separado} y todo complejo
$\mathscr P_\bullet$ de $\mathscr O_X$-módulos cuasicoherentes, acotado
inferiormente, se tiene un isomorfismo canónico de $\partial$-funtores en
$\mathscr P_\bullet$
\begin{env}[6.7.4.1]
\label{III.6.7.4.1}
\[
  \mathfrak{Tor}_n^Y(f,1_Y;\mathscr P_\bullet,\mathscr O_Y)
  \xrightarrow{\sim}
  \mathscr H^{-n}(f,\mathscr P_\bullet)
  \qquad\text{(para todo $n\in\mathbf Z$)}.
  \tag{6.7.4.1}\label{III.6.7.4.1.eq}
\]
\end{env}
En efecto, los métodos de pegado de (6.7.2) reducen inmediatamente al caso en
que $Y$ es afín; entonces, en virtud de (6.2.2), pueden calcularse los dos
miembros de (6.7.4.1) mediante un mismo cubrimiento finito $\mathfrak U$ de
$X$ por abiertos afines y, para el primer miembro, mediante el cubrimiento de
$Y$ formado por el propio $Y$; con las notaciones de (6.6.2), el bicomplejo
$L_{\bullet\bullet}^{(2)}$ se reduce entonces a su término de grados $(0,0)$,
igual a $A$, y la conclusión se deduce de ($0_{\mathrm I}$, 11.7.5). Para una
generalización de este resultado, véase (6.7.7); pero obsérvese que, cuando en
el primer miembro de (6.7.4.1) se reemplaza $\mathscr O_Y$ por un
$\mathscr O_Y$-módulo cuasicoherente \emph{cualquiera} $\mathscr F$, ya no
se tiene en general un isomorfismo con
$\mathscr H^{-n}(f,\mathscr P_\bullet\otimes_{\mathscr O_Y}\mathscr F)$,
aunque, en el cálculo precedente, el bicomplejo
$L_{\bullet\bullet}^{(1)}\otimes_A L_{\bullet\bullet}^{(2)}$ todavía se
identifica con el bicomplejo
$C^\bullet(\mathfrak U,\mathscr P_\bullet\otimes_Y\mathscr F)$.

Por otra parte, se tiene un isomorfismo canónico de
bi-$\partial$-funtores
\begin{env}[6.7.4.2]
\label{III.6.7.4.2}
\[
  \mathfrak{Tor}_\bullet^S
  (1_{X^{(1)}},1_{X^{(2)}};
   \mathscr P_\bullet^{(1)},\mathscr P_\bullet^{(2)})
  \xrightarrow{\sim}
  \mathscr{T}\!or_\bullet^S
  (\mathscr P_\bullet^{(1)},\mathscr P_\bullet^{(2)}).
  \tag{6.7.4.2}\label{III.6.7.4.2.eq}
\]
\end{env}
En efecto, por (6.7.1) y (6.7.2), nos reducimos de nuevo al caso en que $S$ y
los $X^{(i)}$ son afines; en el cálculo del primer miembro de (6.7.4.2), puede
tomarse entonces como
\oldpage[III]{160}
cubrimiento $\mathfrak U^{(i)}$ la familia reducida al
único elemento $X^{(i)}$, de modo que, con las notaciones de (6.6.2),
$L_{\bullet\bullet}^{(i)}$ se reduce a
$\Gamma(X^{(i)},\mathscr P_\bullet^{(i)})$ (considerado como un bicomplejo
cuyos términos de primer grado $\ne0$ son nulos), y la igualdad de los dos
miembros de (6.7.4.2) se deduce de (6.4.1.1) y (6.3.1).
\end{remark}

\begin{proposition}[6.7.5]
\label{III.6.7.5}
Sea $u:\mathscr P_\bullet^{(1)}\to\mathscr Q_\bullet^{(1)}$ un
homomorfismo de complejos acotados inferiormente de
$\mathscr O_{X^{(1)}}$-módulos cuasicoherentes tal que el homomorfismo
\[
  \mathscr H_\bullet(u):
  \mathscr H_\bullet(\mathscr P_\bullet^{(1)})
  \longrightarrow
  \mathscr H_\bullet(\mathscr Q_\bullet^{(1)})
\]
inducido por $u$ sea un isomorfismo. Entonces los homomorfismos
\[
  {}^{(t)}\mathscr E
  (f_1,f_2;\mathscr P_\bullet^{(1)},\mathscr P_\bullet^{(2)})
  \longrightarrow
  {}^{(t)}\mathscr E
  (f_1,f_2;\mathscr Q_\bullet^{(1)},\mathscr P_\bullet^{(2)})
\]
inducidos por $u$ son isomorfismos para $t=a$, $t=b$ y $t=c$.
\end{proposition}

\begin{proof}
La afirmación relativa a la sucesión espectral $(c)$ se deduce de que esta
sucesión es birregular y de que, por hipótesis, el homomorfismo considerado es
un isomorfismo sobre los términos $E_2$ ($0_{\mathrm I}$, 11.1.5). Esto
demuestra ya que
\[
  \mathfrak{Tor}_\bullet^S
  (f_1,f_2;\mathscr P_\bullet^{(1)},\mathscr P_\bullet^{(2)})
  \longrightarrow
  \mathfrak{Tor}_\bullet^S
  (f_1,f_2;\mathscr Q_\bullet^{(1)},\mathscr P_\bullet^{(2)})
\]
es un isomorfismo. Aplicando las relaciones (6.7.4.1) y (6.7.4.2), se ve
primero que los homomorfismos
\[
  \mathscr H^{-n}(f_1,\mathscr P_\bullet^{(1)})
  \longrightarrow
  \mathscr H^{-n}(f_1,\mathscr Q_\bullet^{(1)})
  \quad\text{y}\quad
  \mathscr{T}\!or_n^S
  (\mathscr P_\bullet^{(1)},\mathscr P_\bullet^{(2)})
  \longrightarrow
  \mathscr{T}\!or_n^S
  (\mathscr Q_\bullet^{(1)},\mathscr P_\bullet^{(2)})
\]
inducidos por $u$ son isomorfismos. La afirmación relativa a las sucesiones
$(a)$ y $(b)$ se deduce entonces de que estas sucesiones son birregulares
(6.7.3) y de que los homomorfismos considerados son biyectivos sobre los
términos $E_2$ ($0_{\mathrm I}$, 11.1.5).

Obsérvese, por otra parte, que, si
$u:\mathscr P_\bullet^{(1)}\to\mathscr Q_\bullet^{(1)}$ es una
\emph{equivalencia homotópica}, se deducen de ella isomorfismos canónicos
\[
  {}^{(t)}\mathscr E
  (f_1,f_2;\mathscr P_\bullet^{(1)},\mathscr P_\bullet^{(2)})
  \longrightarrow
  {}^{(t)}\mathscr E
  (f_1,f_2;\mathscr Q_\bullet^{(1)},\mathscr P_\bullet^{(2)})
\]
para las seis sucesiones espectrales. En efecto, si $S$ y los $Y^{(i)}$ son
afines, se deduce de $u$ una equivalencia homotópica de bicomplejos
\[
  C^\bullet(\mathfrak U^{(1)},\mathscr P_\bullet^{(1)})
  \longrightarrow
  C^\bullet(\mathfrak U^{(1)},\mathscr Q_\bullet^{(1)}),
\]
y la proposición se deduce de la teoría general de la hiperhomología
($0_{\mathrm I}$, 11.3.2); el paso al caso general se efectúa por pegado,
utilizando el hecho de que una equivalencia homotópica de complejos induce
una equivalencia homotópica entre resoluciones proyectivas de
Cartan--Eilenberg de dichos complejos (M, XVII, 1.2).
\end{proof}

\begin{proposition}[6.7.6]
\label{III.6.7.6}
Supongamos que el complejo $\mathscr P_\bullet^{(1)}$ o el complejo
$\mathscr P_\bullet^{(2)}$ está formado por módulos $S$-planos (estando ambos
complejos acotados inferiormente). Se tiene entonces un isomorfismo canónico
de bi-$\partial$-funtores
\begin{env}[6.7.6.1]
\label{III.6.7.6.1}
\[
  \mathfrak{Tor}_n^S
  (f_1,f_2;\mathscr P_\bullet^{(1)},\mathscr P_\bullet^{(2)})
  \xrightarrow{\sim}
  \mathscr H^{-n}\bigl(f_1\times_S f_2,
    \mathscr P_\bullet^{(1)}\otimes_S\mathscr P_\bullet^{(2)}\bigr).
  \tag{6.7.6.1}\label{III.6.7.6.1.eq}
\]
\end{env}
\end{proposition}

\begin{proof}
Supongamos primero que $S$, $Y^{(1)}$ e $Y^{(2)}$ son afines, de modo que
estamos en la situación de (6.6.2), cuyas notaciones conservamos. Supongamos,
por ejemplo, que $\mathscr P_\bullet^{(1)}$ está formado por módulos
$S$-planos, y calculemos el hipertor utilizando la observación (6.6.6): es,
por tanto, la homología de
$L_{\bullet\bullet}^{(1)}\otimes_A M_{\bullet\bullet\bullet}^{(2)}$, donde
$M_{\bullet\bullet\bullet}^{(2)}$ es una resolución proyectiva de
Cartan--Eilenberg de $L_{\bullet\bullet}^{(2)}$, en el sentido de
($0_{\mathrm I}$, 11.7.1). Por otra parte, los módulos
$L_{ij}^{(1)}$ son \emph{planos} sobre $A$ en virtud de la hipótesis
(I, 4.15.1); de ($0_{\mathrm I}$, 11.7.5) se deduce entonces un isomorfismo
canónico
\begin{env}[6.7.6.2]
\label{III.6.7.6.2}
\[
  \mathbf{Tor}_\bullet^S
  (\mathfrak U^{(1)},\mathfrak U^{(2)};
   \mathscr P_\bullet^{(1)},\mathscr P_\bullet^{(2)})
  \xrightarrow{\sim}
  \mathscr H_\bullet
  (L_{\bullet\bullet}^{(1)}\otimes_A L_{\bullet\bullet}^{(2)}).
  \tag{6.7.6.2}\label{III.6.7.6.2.eq}
\]
\end{env}
Por otra parte, se tiene un homomorfismo natural de bicomplejos de
$L_{\bullet\bullet}^{(1)}\otimes_A L_{\bullet\bullet}^{(2)}$ en
$C^\bullet(\mathfrak U,\mathscr Q_\bullet)$, donde $\mathfrak U$ es el
cubrimiento de $Z=X^{(1)}\times_S X^{(2)}$ por los abiertos afines
\oldpage[III]{161}
$U_\alpha^{(1)}\times_S U_\beta^{(2)}$ y
$\mathscr Q_\bullet
=\mathscr P_\bullet^{(1)}\otimes_S\mathscr P_\bullet^{(2)}$ (considerado como
complejo simple para el grado total); en efecto, la definición de este
homomorfismo se dio en sustancia durante el cálculo de la sucesión $(b')$ en
(6.6.4), para $q=0$; basta simplemente (conservando las notaciones de (6.6.4))
tener en cuenta que hay, por una parte, un homomorfismo natural del complejo
$N^{\bullet k}$ en el complejo
$C^\bullet(\Phi^{(1)},\Phi^{(2)};\mathscr S_k)$
($0_{\mathrm I}$, 11.8.5), por otra parte, un homomorfismo natural de este
último complejo en el complejo
$P^\bullet(\Phi^{(1)},\Phi^{(2)};\mathscr S_k)$
($0_{\mathrm I}$, 11.8.6), y, por último, un homomorfismo natural de este
último complejo de cocadenas en el subcomplejo de las cocadenas alternadas
($0_{\mathrm I}$, 11.8.7). Además, el homomorfismo de bicomplejos así definido
induce un isomorfismo en homología, como se vio en (6.6.4); componiendo con
(6.7.6.2), se obtiene, por tanto, un isomorfismo
\begin{env}[6.7.6.3]
\label{III.6.7.6.3}
\[
  \mathbf{Tor}_n^S
  (\mathfrak U^{(1)},\mathfrak U^{(2)};
   \mathscr P_\bullet^{(1)},\mathscr P_\bullet^{(2)})
  \xrightarrow{\sim}
  \mathbf H^{-n}
  (\mathfrak U,\mathscr P_\bullet^{(1)}\otimes_S\mathscr P_\bullet^{(2)}).
  \tag{6.7.6.3}\label{III.6.7.6.3.eq}
\]
\end{env}
Hay que demostrar a continuación que el isomorfismo así definido no depende
de los cubrimientos abiertos elegidos (el segundo miembro de (6.7.6.3) es
canónicamente isomorfo a
$\mathbf H^{-n}
(X^{(1)}\times_S X^{(2)},
 \mathscr P_\bullet^{(1)}\otimes_S\mathscr P_\bullet^{(2)})$ por (6.2.2));
esto se hace mediante (6.6.7), observando (con las notaciones de (6.6.7)) que
se tiene un diagrama conmutativo salvo homotopía
\[
\xymatrix@C=18pt@R=28pt{
  C^\bullet(\mathfrak U^{(1)},\mathscr P_\bullet^{(1)})
  \otimes_A
  C^\bullet(\mathfrak U^{(2)},\mathscr P_\bullet^{(2)})
  \ar[r]\ar[d]
  & C^\bullet(\mathfrak U,\mathscr Q_\bullet)\ar[d]
  \\
  C^\bullet(\mathfrak V^{(1)},\mathscr P_\bullet^{(1)})
  \otimes_A
  C^\bullet(\mathfrak V^{(2)},\mathscr P_\bullet^{(2)})
  \ar[r]
  & C^\bullet(\mathfrak V,\mathscr Q_\bullet)
}
\]
donde las flechas horizontales son los homomorfismos definidos anteriormente.
Por último, hay que pasar al caso general por pegado, lo que se hace sin
dificultad como en (6.7.1) y (6.7.2); dejamos los detalles al lector.
\end{proof}

\begin{proposition}[6.7.7]
\label{III.6.7.7}
Supongamos que $\mathscr P_\bullet^{(1)}$ y
$\mathscr P_\bullet^{(2)}$ están acotados inferiormente, y que todos los
módulos $\mathscr H^{-n}(f_1,\mathscr P_\bullet^{(1)})$ o todos los módulos
$\mathscr H^{-n}(f_2,\mathscr P_\bullet^{(2)})$ son $S$-planos. Se tiene
entonces un isomorfismo canónico de bi-$\partial$-funtores ($n$ recorre
$\mathbf Z$)
\begin{env}[6.7.7.1]
\label{III.6.7.7.1}
\[
  \mathfrak{Tor}_n^S
  (f_1,f_2;\mathscr P_\bullet^{(1)},\mathscr P_\bullet^{(2)})
  \xrightarrow{\sim}
  \bigoplus_{q_1+q_2=n}
  \mathscr H^{-q_1}(f_1,\mathscr P_\bullet^{(1)})
  \otimes_S
  \mathscr H^{-q_2}(f_2,\mathscr P_\bullet^{(2)}).
  \tag{6.7.7.1}\label{III.6.7.7.1.eq}
\]
\end{env}
\end{proposition}

\begin{proof}
En efecto, en vista de (6.5.8), la sucesión espectral $(a)$ de (6.7.3) es
degenerada, y la proposición se deduce inmediatamente de
($0_{\mathrm I}$, 11.1.6), pues esta sucesión es birregular (6.7.3).
\end{proof}

\begin{theorem}[6.7.8]
\label{III.6.7.8}
Supongamos que se cumplen las condiciones siguientes:
\begin{enumerate}
  \item[{$1^\circ$}]
    Los complejos $\mathscr P_\bullet^{(1)}$ y
    $\mathscr P_\bullet^{(2)}$ están acotados inferiormente.
  \item[{$2^\circ$}]
    El complejo $\mathscr P_\bullet^{(1)}$ o el complejo
    $\mathscr P_\bullet^{(2)}$ está formado por módulos $S$-planos.

  \item[{$3^\circ$}]
    \oldpage[III]{162}
    Todos los módulos $\mathscr H^{-n}(f_1,\mathscr P_\bullet^{(1)})$ o todos
    los módulos $\mathscr H^{-n}(f_2,\mathscr P_\bullet^{(2)})$ son
    $S$-planos.
\end{enumerate}
Se tiene entonces un isomorfismo canónico de bi-$\partial$-funtores
($n$ recorre $\mathbf Z$)
\begin{env}[6.7.8.1]
\label{III.6.7.8.1}
\[
  \mathscr H^n\bigl(f_1\times_S f_2,
    \mathscr P_\bullet^{(1)}\otimes_S\mathscr P_\bullet^{(2)}\bigr)
  \xrightarrow{\sim}
  \bigoplus_{n_1+n_2=n}
  \mathscr H^{n_1}(f_1,\mathscr P_\bullet^{(1)})
  \otimes_S
  \mathscr H^{n_2}(f_2,\mathscr P_\bullet^{(2)}).
  \tag{6.7.8.1}\label{III.6.7.8.1.eq}
\]
\end{env}
(«fórmula de Künneth»).
\end{theorem}

\begin{proof}
Esto se deduce de (6.7.6) y (6.7.7).

Cuando $S$, $Y^{(1)}$ e $Y^{(2)}$ son afines, el isomorfismo recíproco de
(6.7.8.1) se deduce (con las notaciones de (6.7.6)) del homomorfismo de
bicomplejos
\[
  C^\bullet(\mathfrak U^{(1)},\mathscr P_\bullet^{(1)})
  \otimes_A
  C^\bullet(\mathfrak U^{(2)},\mathscr P_\bullet^{(2)})
  \longrightarrow
  C^\bullet(\mathfrak U,\mathscr Q_\bullet)
\]
mediante el procedimiento definido en (G, I, 2.7), como se deduce de
(G, I, 5.5).
\end{proof}

\begin{proposition}[6.7.9]
\label{III.6.7.9}
Supongamos que se cumplen las tres condiciones siguientes:
\begin{enumerate}
  \item[{$1^\circ$}]
    $S$, $Y^{(1)}$ e $Y^{(2)}$ son localmente noetherianos, $f_1$ y $f_2$
    son propios, y $Y^{(1)}$ o $Y^{(2)}$ es de tipo finito sobre $S$.
  \item[{$2^\circ$}]
    $\mathscr P_\bullet^{(1)}$ y $\mathscr P_\bullet^{(2)}$ están acotados
    inferiormente.
  \item[{$3^\circ$}]
    Para todo $n\in\mathbf Z$,
    $\mathscr H_n(\mathscr P_\bullet^{(i)})$ es un módulo coherente
    ($i=1,2$).
\end{enumerate}
En estas condiciones,
$\mathfrak{Tor}_n^S
(f_1,f_2;\mathscr P_\bullet^{(1)},\mathscr P_\bullet^{(2)})$ es un
$\mathscr O_Y$-módulo coherente (con
$Y=Y^{(1)}\times_S Y^{(2)}$).
\end{proposition}

\begin{proof}
De (6.5.13) se deduce que los hipertor locales
$\mathscr{T}\!or_n^S(\mathscr P_\bullet^{(1)},\mathscr P_\bullet^{(2)})$ son
$\mathscr O_X$-módulos coherentes
($X=X^{(1)}\times_S X^{(2)}$ es localmente noetheriano, pues uno de los
$X^{(i)}$ es, por hipótesis, de tipo finito sobre $S$ (I, 6.3.4 y 6.3.8)).
Como $Y$ es localmente noetheriano y $f_1\times_S f_2$ es propio
(II, 5.4.2), de (6.2.5) se deduce que los términos
${}^{(b)}\mathscr E_{pq}^2$ de (6.7.3) son $\mathscr O_Y$-módulos coherentes.
Como todas las sucesiones espectrales de (6.7.3) son birregulares en virtud
de la hipótesis $2^\circ$, se concluye por ($0_{\mathrm I}$, 11.1.8).
\end{proof}

\begin{env}[6.7.10]
\label{III.6.7.10}
Sean ahora $Y'^{(i)}$ dos $S$-preesquemas,
$v_i:Y'^{(i)}\to Y^{(i)}$ dos $S$-morfismos ($i=1,2$) y
$v=v_1\times_S v_2$ su producto, que es un $S$-morfismo $Y'\to Y$, donde se
pone $Y'=Y'^{(1)}\times_S Y'^{(2)}$. Consideremos, por otra parte, para
$i=1,2$, un $S$-preesquema $X'^{(i)}$ y dos $S$-morfismos
$u_i:X'^{(i)}\to X^{(i)}$, $f'_i:X'^{(i)}\to Y'^{(i)}$, de modo que los
diagramas
\begin{env}[6.7.10.1]
\label{III.6.7.10.1}
\[
\xymatrix@C=36pt@R=28pt{
  X'^{(i)}\ar[r]^{u_i}\ar[d]_{f'_i}
  & X^{(i)}\ar[d]^{f_i}
  \\
  Y'^{(i)}\ar[r]_{v_i}
  & Y^{(i)}
}
\tag{6.7.10.1}\label{III.6.7.10.1.eq}
\]
\end{env}
sean conmutativos, siendo los morfismos $f'_i$ \emph{separados y
cuasicompactos}. Se tienen entonces homomorfismos canónicos
$\mathscr O_{Y'}$-lineales de funtores espectrales
\begin{env}[6.7.10.2]
\label{III.6.7.10.2}
\[
  v^*\bigl({}^{(t)}\mathscr E
  (f_1,f_2;\mathscr P_\bullet^{(1)},\mathscr P_\bullet^{(2)})\bigr)
  \longrightarrow
  {}^{(t)}\mathscr E
  (f'_1,f'_2;u_1^*(\mathscr P_\bullet^{(1)}),
                 u_2^*(\mathscr P_\bullet^{(2)}))
  \tag{6.7.10.2}\label{III.6.7.10.2.eq}
\]
\end{env}
para $t=a$, $a'$, $b$, $b'$, $c$, $d$. Para definirlos, supongamos primero
que $S=\operatorname{Spec}(A)$,
$Y^{(i)}=\operatorname{Spec}(B_i)$ e
$Y'^{(i)}=\operatorname{Spec}(B'_i)$ son afines; entonces $X^{(i)}$ y
$X'^{(i)}$ son esquemas cuasicompactos. Para calcular las sucesiones
espectrales
${}^{(t)}\mathscr E
(f_1,f_2;\mathscr P_\bullet^{(1)},\mathscr P_\bullet^{(2)})$, consideramos,
como en (6.6.1), cubrimientos finitos $\mathfrak U^{(i)}$ por abiertos afines
de $X^{(i)}$ ($i=1,2$); para calcular
${}^{(t)}\mathscr E
(f'_1,f'_2;u_1^*(\mathscr P_\bullet^{(1)}),
u_2^*(\mathscr P_\bullet^{(2)}))$, consideramos cubrimientos finitos
$\mathfrak U'^{(i)}$

\oldpage[III]{163}
de $X'^{(i)}$ por abiertos afines, \emph{más finos}, respectivamente, que los
cubrimientos $u_i^{-1}(\mathfrak U^{(i)})$ ($i=1,2$). Está claro que el
bicomplejo
$C^\bullet(\mathfrak U^{(i)},\mathscr P_\bullet^{(i)})
=L_{\bullet\bullet}^{(i)}$ puede considerarse canónicamente como un
subbicomplejo de
$C^\bullet(u_i^{-1}(\mathfrak U^{(i)}),u_i^*(\mathscr P_\bullet^{(i)}))$
($0_{\mathrm I}$, 4.4.3.2); además, eligiendo un morfismo simplicial
(G, II, 5.7) de $\mathfrak U'^{(i)}$ en
$u_i^{-1}(\mathfrak U^{(i)})$, se define un homomorfismo de bicomplejos
\[
  C^\bullet(u_i^{-1}(\mathfrak U^{(i)}),u_i^*(\mathscr P_\bullet^{(i)}))
  \longrightarrow
  C^\bullet(\mathfrak U'^{(i)},u_i^*(\mathscr P_\bullet^{(i)}))
\]
del cual se obtiene, por composición, un homomorfismo de bicomplejos
\[
  L_{\bullet\bullet}^{(i)}
  \longrightarrow
  L_{\bullet\bullet}^{\prime(i)}
  =C^\bullet(\mathfrak U'^{(i)},u_i^*(\mathscr P_\bullet^{(i)})).
\]
Además, este homomorfismo se sustituye por uno homotópico cuando se cambia el
morfismo simplicial (G, II, 5.7.1); así se obtiene un homomorfismo bien
definido de funtores espectrales:
\begin{env}[6.7.10.3]
\label{III.6.7.10.3}
\[
  {}^{(t)}\mathscr E
  (\mathfrak U^{(1)},\mathfrak U^{(2)};
   \mathscr P_\bullet^{(1)},\mathscr P_\bullet^{(2)})
  \longrightarrow
  {}^{(t)}\mathscr E
  (\mathfrak U'^{(1)},\mathfrak U'^{(2)};
   u_1^*(\mathscr P_\bullet^{(1)}),u_2^*(\mathscr P_\bullet^{(2)})).
  \tag{6.7.10.3}\label{III.6.7.10.3.eq}
\]
\end{env}
Se verifica inmediatamente que, si $\mathfrak V^{(i)}$ es un cubrimiento
afín finito de $X^{(i)}$ más fino que $\mathfrak U^{(i)}$, y
$\mathfrak V'^{(i)}$ un cubrimiento afín finito de $X'^{(i)}$, más fino que
$u_i^{-1}(\mathfrak V^{(i)})$ y que $\mathfrak U'^{(i)}$, el diagrama
\[
\xymatrix@C=34pt@R=30pt{
  C^\bullet(\mathfrak U^{(i)},\mathscr P_\bullet^{(i)})
  \ar[r]\ar[d]
  & C^\bullet(\mathfrak V^{(i)},\mathscr P_\bullet^{(i)})\ar[d]
  \\
  C^\bullet(\mathfrak U'^{(i)},u_i^*(\mathscr P_\bullet^{(i)}))
  \ar[r]
  & C^\bullet(\mathfrak V'^{(i)},u_i^*(\mathscr P_\bullet^{(i)}))
}
\]
es conmutativo, lo que implica que el homomorfismo (6.7.10.3) no depende
esencialmente de los cubrimientos $\mathfrak U^{(i)}$ y
$\mathfrak U'^{(i)}$ considerados. Por tanto, se ha definido en realidad un
homomorfismo de $A$-módulos
\begin{env}[6.7.10.4]
\label{III.6.7.10.4}
\[
  {}^{(t)}E
  (X^{(1)},X^{(2)};\mathscr P_\bullet^{(1)},\mathscr P_\bullet^{(2)})
  \longrightarrow
  {}^{(t)}E
  (X'^{(1)},X'^{(2)};
   u_1^*(\mathscr P_\bullet^{(1)}),u_2^*(\mathscr P_\bullet^{(2)})).
  \tag{6.7.10.4}\label{III.6.7.10.4.eq}
\]
\end{env}
pero, por la definición de los $u_i^*(\mathscr P_\bullet^{(i)})$ y en virtud
de la conmutatividad de (6.7.10.1), está claro que este homomorfismo es también
un homomorfismo de $(B_1\otimes_A B_2)$-módulos; como el segundo miembro de
(6.7.10.4) está formado por $(B'_1\otimes_A B'_2)$-módulos, se deduce
canónicamente de (6.7.10.4) un homomorfismo de
$(B'_1\otimes_A B'_2)$-módulos
\begin{env}[6.7.10.5]
\label{III.6.7.10.5}
\[
  {}^{(t)}E
  (X^{(1)},X^{(2)};\mathscr P_\bullet^{(1)},\mathscr P_\bullet^{(2)})
  \otimes_{B_1\otimes_A B_2}(B'_1\otimes_A B'_2)
  \longrightarrow
  {}^{(t)}E
  (X'^{(1)},X'^{(2)};
   u_1^*(\mathscr P_\bullet^{(1)}),u_2^*(\mathscr P_\bullet^{(2)}))
  \tag{6.7.10.5}\label{III.6.7.10.5.eq}
\]
\end{env}
lo que, teniendo en cuenta (I, 1.6.5), no es sino el homomorfismo buscado
(6.7.10.2) en el caso particular considerado.

Queda pasar al caso general siguiendo los procedimientos de pegado de (6.7.1)
y (6.7.2); el
segundo paso es inmediato; por lo que respecta al primero, se consideran, como
en (6.7.1), elementos $g_i\in B_i$ y sus imágenes $g'_i\in B'_i$, el producto
tensorial $g=g_1\otimes g_2$ en $B=B_1\otimes_A B_2$ y su imagen
$g'=g'_1\otimes g'_2$ en $B'=B'_1\otimes_A B'_2$, y todo se reduce a utilizar

\oldpage[III]{164}
el isomorfismo canónico
\[
  (M\otimes_B B')_{g'}
  \xrightarrow{\sim}
  M_g\otimes_{B_g}B'_{g'}
\]
($0_{\mathrm I}$, 1.5.4); dejamos los detalles al lector.
\end{env}

\begin{env}[6.7.11]
\label{III.6.7.11}
La teoría de los hipertor globales, desarrollada anteriormente para dos
$S$-morfismos $X^{(i)}\to Y^{(i)}$ y dos complejos
$\mathscr P_\bullet^{(i)}$ de módulos cuasicoherentes acotados inferiormente,
se extiende inmediatamente al caso general siguiente: se tiene un preesquema
$S$, una familia finita de $S$-preesquemas $Y^{(i)}$ ($i\in I$), una familia
finita de $S$-morfismos separados y cuasicompactos
$f_i:X^{(i)}\to Y^{(i)}$, y, para cada $i$, un complejo acotado inferiormente
de $\mathscr O_{X^{(i)}}$-módulos cuasicoherentes
$\mathscr P_\bullet^{(i)}$. Si $Y$ es el producto de los $S$-preesquemas
$Y^{(i)}$, se define entonces, para cada entero $n\in\mathbf Z$, un
$\mathscr O_Y$-módulo cuasicoherente
$\mathfrak{Tor}_n^S((f_i)_{i\in I};
(\mathscr P_\bullet^{(i)})_{i\in I})$;
estos módulos forman un $\partial$-funtor covariante en cada uno de los
complejos $\mathscr P_\bullet^{(i)}$; además, este funtor es el término límite
común de seis funtores espectrales
${}^{(t)}\mathscr E((f_i)_{i\in I};(\mathscr P_\bullet^{(i)})_{i\in I})$.
Dejamos al lector la tarea de repetir, para este caso general, las definiciones
y los razonamientos hechos anteriormente para $I=\{1,2\}$. Obsérvese
simplemente que, cuando $I$ se reduce a un solo elemento, se recupera la
hipercohomología $\mathscr H^\bullet(f,\mathscr P_\bullet)$ definida en
(6.2.7) (como ya se observó en (6.7.4)). Cuando $I$ es el intervalo
$1\leq i\leq m$ de $\mathbf N$, escribiremos
\phantomsection\label{III.6.7.11.2}
\[
  \mathfrak{Tor}_n^S
  (f_1,\ldots,f_m;
   \mathscr P_\bullet^{(1)},\ldots,\mathscr P_\bullet^{(m)})
  \quad\text{para}\quad
  \mathfrak{Tor}_n^S
  ((f_i)_{i\in I};(\mathscr P_\bullet^{(i)})_{i\in I}).
\]
\end{env}

\begin{proposition}[6.7.12]
\label{III.6.7.12}
Con las notaciones de (6.7.11), sea $J$ un subconjunto de $I$ tal que, para
$i\in I-J$, se tenga $X^{(i)}=Y^{(i)}=S$, que $f_i$ sea la identidad y que
$\mathscr P_\bullet^{(i)}$ sea el complejo reducido al término de grado $0$
igual a $\mathscr O_S$. Hay entonces un isomorfismo canónico de
$\partial$-funtores
\begin{env}[6.7.12.1]
\label{III.6.7.12.1}
\[
  \mathfrak{Tor}_\bullet^S
  ((f_i)_{i\in I};(\mathscr P_\bullet^{(i)})_{i\in I})
  \xrightarrow{\sim}
  \mathfrak{Tor}_\bullet^S
  ((f_i)_{i\in J};(\mathscr P_\bullet^{(i)})_{i\in J}).
  \tag{6.7.12.1}\label{III.6.7.12.1.eq}
\]
\end{env}
\end{proposition}

\begin{proof}
Podemos limitarnos a definir este isomorfismo cuando $S$ y los $Y^{(i)}$ son
afines, pues el pegado se efectúa como de costumbre. Para $i\in I-J$, puede
tomarse el cubrimiento $\mathfrak U^{(i)}$ formado por el único conjunto $S$,
y entonces
$L_{\bullet\bullet}^{(i)}=C^\bullet(\mathfrak U^{(i)},
\mathscr P_\bullet^{(i)})$ se reduce a su único término de grados $(0,0)$,
igual a $\Gamma(S,\mathscr O_S)=A(S)$; el isomorfismo (6.7.12.1) es entonces
evidente.
\end{proof}

\begin{remark}[6.7.13]
\label{III.6.7.13}
Con las notaciones de (6.7.3), consideremos el $S$-isomorfismo canónico
$Y^{(1)}\times_S Y^{(2)}\to Y^{(2)}\times_S Y^{(1)}$ (I, 3.3.5); entonces la
imagen por este isomorfismo de
$\mathfrak{Tor}_\bullet^S(f_1,f_2;
\mathscr P_\bullet^{(1)},\mathscr P_\bullet^{(2)})$ es
$\mathfrak{Tor}_\bullet^S(f_2,f_1;
\mathscr P_\bullet^{(2)},\mathscr P_\bullet^{(1)})$; como la cuestión es
local, nos reducimos al caso considerado en (6.6.2), y, si se designa por
$M_{\bullet\bullet\bullet}^{(i)}$ una resolución proyectiva de
Cartan--Eilenberg de $L_{\bullet\bullet}^{(i)}$ ($i=1,2$), el isomorfismo
considerado transforma
$M_{\bullet\bullet\bullet}^{(1)}\otimes_A
M_{\bullet\bullet\bullet}^{(2)}$ en
$M_{\bullet\bullet\bullet}^{(2)}\otimes_A
M_{\bullet\bullet\bullet}^{(1)}$, de donde se deduce la afirmación al
considerar la homología de los complejos simples asociados a estos
tricomplejos.
\end{remark}
